\documentclass[11pt]{article}
\usepackage[T1]{fontenc}
\usepackage{amsmath}
\usepackage{booktabs}
\usepackage[hidelinks]{hyperref}
\usepackage[a4paper,margin=25mm]{geometry}

\newcommand{\labid}{golden-basic-pdflatex}

\begin{document}
\title{Golden Matrix Baseline (pdfLaTeX)}
\author{Company Agent LaTeX Toolchain Lab}
\date{2026-10-26}
\maketitle

\section{Purpose}
This baseline project exercises the pdfLaTeX path: sections, cross references,
an equation, and a table. The golden matrix compiles it without shell escape
and compares the page count against the recorded Overleaf reference when one
is supplied.

\section{Method}\label{sec:method}
The loss function is defined in Equation~\eqref{eq:loss}. This section uses
Table~\ref{tab:devices} to document the devices referenced by the lab.

\begin{equation}\label{eq:loss}
  \mathcal{L} = \sum_{i=1}^{N} \bigl(f_{\theta}(x_i) - y_i\bigr)^2
\end{equation}

\begin{table}[ht]
  \centering
  \caption{Devices referenced by the golden lab}
  \label{tab:devices}
  \begin{tabular}{llr}
    \toprule
    Device & Region & Qubits \\
    \midrule
    iontrap-a & Hefei & 5 \\
    simulator-b & Local & 5 \\
    \bottomrule
  \end{tabular}
\end{table}

\section{Expected Result}
The compiled PDF has at least one page and contains the baseline identifier
\texttt{\labid}. Page geometry stays A4 with 25\,mm margins.

\end{document}
